% ----------------------
\subsection{Section 7 (April 1829)}
% ----------------------
	\label{DC7}
	
	% -----
	\subsubsection{Historical Background}
	% -----
		\label{DC7-HB}
		
		% --
		\paragraph{Commentary on this document as a ``translation''}
		% --
		
			HDDC says that the manuscript versions of this item do not describe it as a ``translation''; they instead describe it as a ``revelation.'' This is true for the account in the \hyperref[EMS-1828-1829-JS-APR-1829-C]{JSP}, which comes from Revelation Book 1, but HDDC (page 177) also mentions two other manuscript versions of this section, both of which (apparently) describe it as a ``revelation.'' I haven't yet been able to get access to these manuscript versions, but I will email the CHL sometime about it.
			
			The important thing here is that, as far as I can tell based on the evidence I have seen, this document was called a ``translation'' for the first time in the BC version in 1833; before that, if they called it a ``translation,'' no one had made a note of that.
			
			Although, if you see \hyperref[DC8]{DC 8} as coming before this revelation, then there is stuff in 8 (including in the original version) that could be some slight evidence that they were thinking about section 7 as a translation from the beginning. In 8 I'm thinking of \hyperref[DC8-11]{verse 11} in particular.
			
		% --
		\paragraph{The ``FGW'' version}
		% --
		
			This is a copy in the handwriting of Frederick G. Williams, dated by the JSP researchers to around 1836. View it on the JSP website \href{https://www.josephsmithpapers.org/paper-summary/account-of-john-april-1829-c-frederick-g-williams-copy-dc-7/1}{here}; see the whole document \href{https://drive.google.com/file/d/1BfD8CjmUlMydDjPZAj9Y5HwFHz2WgTe8/view?usp=sharing}{here}. The JSP source note is, ``Account of John, [Harmony Township, Susquehanna Co., PA, Apr. 1829]. Version copied [ca. 1836]; handwriting of Frederick G. Williams; one page; Revelations Collection, CHL.'' 
			
			HDDC mentions a manuscript version in the handwriting of Williams that I believe is this version, on page 177. He dates it to ``Prior to March 27, 1836,'' because of some additional material written on it regarding the Kirtland Temple. This additional material is discussed in some typewritten letters attached to the document in the link to the entire thing above.
		
		% --
		\paragraph{JSP}
		% --
		
			``In April 1829, JS dictated the following revelation, which in its first publication was described as the translation of an ancient parchment written by the apostle John.%
				\footnote{%
					% \label{}%
					See my footnote to \hyperref[DC7-1]{verse 1} below.
					}
			Ancient writings besides the Book of Mormon began to interest JS and Oliver Cowdery not long after they started their translation of the gold plates. The Book of Mormon manuscript itself mentioned several ancient texts, and additionally, JS had dictated a revelation promising Cowdery the privilege, if he so desired, of translating `records which contain much of my gospel, which have been kept back because of the wickedness of the people.' Then, as JS and Cowdery continued the Book of Mormon translation, a `difference of opinion' arose between them regarding a question left unanswered in the New Testament: whether `John the Apostle...died, or whether he continued' on earth until the second coming of Christ. The source of this disagreement was the final chapter of the Gospel of John, in which Jesus prophesied of the apostle Peter's death. When Peter asked what would happen to his fellow apostle John, Jesus responded, `If I will that he should tarry till I come, what is that to thee?' Questions about the fate of John were common in JS's time. For example, Adam Clarke, a noted Bible commentator, wrote, `For nearly eighteen hundred years, the greatest men in the world have been puzzled with this passage [\hyperref[JOHN-21-22]{John 21:22}].' JS and Cowdery's discussion of this issue possibly arose when they encountered a passage in the translation of the plates describing the biblical prophet Moses and the Book of Mormon prophet Alma as having been `taken up by the spirit, or buried by the hand of the Lord.'
			
			``JS's history reports that he and Cowdery `mutually agreed to settle it [their question] by the Urim and Thummin, and the following is the word which we received.' As noted, this revelation was said to be `translated from parchment, written and hid up' by John himself, and the text begins in the first person, with John stating, `And the Lord said unto me,' followed by an account in which Jesus declares the respective fates of John and Peter.''
			
		% --
		\paragraph{HDDC}
		% --
		
			``The \textit{Times and Seasons} is the earliest account that indicates this revelation was a translation of an ancient parchment.%
				\footnote{%
					% \label{}%
					This is wrong; see my footnote to \hyperref[DC7-1]{verse 1}.
					}
			Joseph Smith was the editor of the \textit{Times and Seasons} when this revelation was published as part of his history; therefore, this additional item concerning the revelation is an authentic description of its origin. The earlier versions, all manuscripts, contain no mention of the parchment'' (176).
		
	% -----
	\subsubsection{Versions}
	% -----
		\label{DC8-V}
	
		\begin{tabular}{p{1in}p{2in}p{1in}p{2in}}
			JSP		& \hyperref[EMS-1828-1829-JS-APR-1829-C]{April 1829}		& BC	& Chapter 6, 18 \\
			DC 1835	& Section 33, 160--161										& TS	& 3:18 (15 July 1842), 853 \\
			MS		& 3:8 (December 1842), 134									& DC 1844 & Section 33, 236--237 \\
		\end{tabular}
		
		\medskip
		
		See the \href{https://drive.google.com/file/d/1goAvNz8jS4R8slYfMG_db55Ty2z_vmgK/view?usp=sharing}{sections comparison} document for more information about the version history.

	% -----
	\subsubsection{Section 7:1} % *DC7-1 
	% -----
		\label{DC7-1}

			\footnote{%
				% \label{}%
				The JSP version has the heading, ``A Revelation to Joseph \& Oliver [Cowdery] concerning John the Beloved Deciple who leaned on his Saveiours breast given \sout{on} in Harmony Susquehannah County Pennsylvania.'' BC has, ``1 \textit{A Revelation given to Joseph and Oliver [Cowdery], in Harmony, Pennsylvania, April, 1829, when they desired to know whether John, the beloved disciple, tarried on earth. Translated from parchment, written and hid up by himself.}'' HDDC makes a somewhat glaring error about this, which is pretty surprising---on 176 it says ``the \textit{Times and Seasons} is the earliest account that indicates this revelation was a translation of an ancient parchment.'' But this is obviously wrong. DC 1835 has, ``\textit{A Revelation given to Joseph Smith, jr. and Oliver Cowdery, in Harmony, Pennsylvania, April, 1829, when they desired to know whether John, the beloved disciple, tarried on earth.--- Translated from parchment, written and hid up by himself.}'' The FGW version, dated to around 1836, has ``A revelation concerning John the beloved deciple.'' DC 1844 has, ``\textit{Revelation given to Joseph Smith, jr. and Oliver Cowdery, in Harmony, Pennsylvania, April, 1829, when they desired to know whether John the beloved disciple, tarried on earth. Translated from parchment, written and hid up by himself.}''
				}%
		AND the Lord said unto me: John, my beloved, what desirest thou?  For if you shall ask what you will, it shall be granted unto you.

		\begin{framed}
		\scriptsize
			\begin{compactdesc}
				\item[JSP] And the Lord said unto me. John my Beloved what desiredst thou
				\item[BC] AND the Lord said unto me, John my beloved, what desirest thou? 
				\item[DC 1835] 1 And the Lord said unto me, John, my beloved, what desirest thou? For if ye shall ask, what you will, it shall be granted unto you. 
				\item[FGW] And the Lord said unto me John my beloved, what desirest thou 
				\item[DC 1844] 1 And the Lord said unto me, John, my beloved, what desirest thou? For if ye shall ask, what you will, it shall be granted unto you.
			\end{compactdesc}
		\end{framed}

	% -----
	\subsubsection{Section 7:2} % *DC7-2 
	% -----
		\label{DC7-2}

		And I said unto him: Lord, give unto me power over death,%
			\footnote{%
				% \label{}%
				Note that this part about the ``power over death'' is not in the manuscript version nor in the BC. It apparently first appears in the 1835 DC. So you have to think, what happens between 1833 and 1835 such that JS would include something about the ``over death'' part in particular? Not sure how relevant it is, but Oliver Cowdery gave a blessing to Parley Pratt, when Pratt was being ordained one of the twelve by Cowdery, JS, and David Whitmer, on \hyperref[EMS-1835-JS-21-FEB-1835]{21 February 1835}, which includes the remark, ``Let sickness and death have no power over him.''
				}
		that I may live and bring souls unto thee.

		\begin{framed}
		\scriptsize
			\begin{compactdesc}
				\item[JSP] \& I said Lord give unto me power that I may bring souls unto thee 
				\item[BC] and I said Lord, give unto me power that I may bring souls unto thee.---
				\item[DC 1835] And I said unto him, Lord, give unto me power over death, that I may live and bring souls unto thee.
				\item[FGW] And I said Lord give unto me power that I may bring souls unto thee, 
				\item[DC 1844] And I said unto him, Lord, give unto me power over death, that I may live and bring souls unto thee.
			\end{compactdesc}
		\end{framed}

	% -----
	\subsubsection{Section 7:3} % *DC7-3 
	% -----
		\label{DC7-3}

		And the Lord said unto me: Verily, verily, I say unto thee, because thou desirest this thou shalt tarry until I come in my glory, and shalt prophesy before nations, kindreds, tongues and people.

		\begin{framed}
		\scriptsize
			\begin{compactdesc}
				\item[JSP] \& the Lord said unto me Veriley Verily I say unto thee because thou desiredst this thou shalt tarry until I come in my glory
				\item[BC] And the Lord said unto me: Verily, verily I say unto thee, because thou desiredst this, thou shalt tarry till I come in my glory:
				\item[DC 1835] And the Lord said unto me, Verily, verily, I say unto thee, because thou desiredst this thou shalt tarry until I come in my glory, and shall prophesy before nations, kindreds, tongues and people.
				\item[FGW] and the Lord said unto me verily I say unto thee because thou desirest this; thou shalt tarry till I come in my glory
				\item[DC 1844] And the Lord said unto me, Verily, verily I say unto thee, because thou desirest this thou shalt tarry until I come in my glory, and shall prophecy before nations, kindreds, tongues, and people.
			\end{compactdesc}
		\end{framed}

	% -----
	\subsubsection{Section 7:4} % *DC7-4 
	% -----
		\label{DC7-4}

		And for this cause the Lord said unto Peter: If I will that he tarry till I come, what is that to thee?  For he desired of me that he might bring souls unto me, but thou desiredst that thou mightest speedily come unto me in my kingdom.

		\begin{framed}
		\scriptsize
			\begin{compactdesc}
				\item[JSP] \& for this cause the Lord said unto Peter if I will that he tarry till I come what is that to thee For he desiredst of me that he might bring souls unto me but thou desiredst that thou mightest come unto me in my kingdom
				\item[BC] 2 And for this cause, the Lord said unto Peter:--- If I will that he tarry till I come, what is that to thee? for he desiredst of me that he might bring souls unto me: but thou desiredst that thou might speedily come unto me in my kingdom:
				\item[DC 1835] 2 And for this cause the Lord said unto Peter, If I will that he tarry till I come, what is that to thee? For he desiredst of me that he might bring souls unto me; but thou desiredst that thou might speedily come unto me in my kingdom. 
				\item[FGW] and for this cause the Lord said unto Peter if I will that he tarry till I come, what is that to thee for he desirest of me that he might bring souls unto me, but thou desirest that thou might speedily come unto me in my kingdom 
				\item[DC 1844] 2 And for this cause the Lord said unto Peter, If I will that he tarry till I come, what is that to thee? For he desiredst of me that he might bring souls unto me; but thou desiredst that thou might speedily come unto me in my kingdom.
			\end{compactdesc}
		\end{framed}

	% -----
	\subsubsection{Section 7:5} % *DC7-5 
	% -----
		\label{DC7-5}

		I say unto thee, Peter, this was a good desire; but my beloved has desired that he might do more, or a greater work yet among men than what he has before done.

		\begin{framed}
		\scriptsize
			\begin{compactdesc}
				\item[JSP] I say unto thee Peter this was a good desire but my beloved hath undertaken a greater work
				\item[BC] I say unto thee, Peter, this was a good desire, but my beloved has undertaken a greater work.
				\item[DC 1835] I say unto thee, Peter, this was a good desire, but my beloved has desired that he might do more, or a greater work, yet among men than what he has before done;
				\item[FGW] I say unto thee Peter this was a good desire, But my beloved hath undertaken a greater work; 
				\item[DC 1844] I say unto thee, Peter, this was a good desire, but my beloved has desired that he might do more, or a greater work, yet among men than what he has before done, 
			\end{compactdesc}
		\end{framed}

	% -----
	\subsubsection{Section 7:6} % *DC7-6 
	% -----
		\label{DC7-6}

		Yea, he has undertaken a greater work; therefore I will make him as flaming fire%
			\footnote{%
				% \label{}%
				See \hyperref[PSALMS-104-4]{Psalms 104:4} for ``flaming fire'' in this context. See also \hyperref[ISAIAH-4-5]{Isaiah 4:5}. This is not in the JSP or BC, of course.
				}
		and a ministering angel; he shall minister for those who shall be heirs of salvation who dwell on the earth.

		\begin{framed}
		\scriptsize
			\begin{compactdesc}
				\item[JSP] ---
				\item[BC] ---
				\item[DC 1835] yea, he has undertaken a greater work; therefore, I will make him as flaming fire and a ministering angel: he shall minister for those who shall be heirs of salvation who dwell on the earth;
				\item[FGW] ---
				\item[DC 1844] yea, he has undertaken a greater work; therefore, I will make him as flaming fire and a ministering angel: he shall minister for those who shall be heirs of salvation who dwell on the earth;
			\end{compactdesc}
		\end{framed}

	% -----
	\subsubsection{Section 7:7} % *DC7-7 
	% -----
		\label{DC7-7}

		And I will make thee to minister for him and for thy brother James; and unto you three I will give this power and the keys of this ministry until I come.%
			\footnote{%
				% \label{}%
				This looks like endowment stuff to me. 
				}

		\begin{framed}
		\scriptsize
			\begin{compactdesc}
				\item[JSP] ---
				\item[BC] ---
				\item[DC 1835] and I will make thee to minister for him and for thy brother James: and unto you three I will give this power and the keys of this ministry until I come.
				\item[FGW] ---
				\item[DC 1844] and I will make thee to minister for him and thy brother James: and unto you three I will give this power and the keys of the ministry until I come.
			\end{compactdesc}
		\end{framed}

	% -----
	\subsubsection{Section 7:8} % *DC7-8 
	% -----
		\label{DC7-8}

		Verily I say unto you, ye shall both have according to your desires, for ye both joy in that which ye have desired.

		\begin{framed}
		\scriptsize
			\begin{compactdesc}
				\item[JSP] Verily I say unto you ye shall both have according to your desires for ye both Joy in that which ye desired
				\item[BC] 3 Verily I say unto you, ye shall both have according to your desires, for ye both joy in that which ye have desired.
				\item[DC 1835] 3 Verily I say unto you, ye shall both have according to your desires, for ye both joy in that which ye have desired.
				\item[FGW] verily I say unto you ye <shall> both have according to your desires for ye hath joy in th[at] which ye have desired \&c \&c \&c,------
				\item[DC 1844] 3 Verily I say unto you, ye shall both have according to your desires, for ye both joy in that which ye have desired.
			\end{compactdesc}
		\end{framed}